\section{Review Protocol v0.2}

This appendix reproduces the protocol that governed all screening, extraction,
quality assurance, and second-reader verification steps for this review. The
protocol was finalised as Version~0.2 and followed without deviation. The
original file (\texttt{protocol\_v0.2.md}) is archived alongside other
repository artefacts.

\subsection{Purpose and Scope}

This protocol defines the procedures for identifying, screening, extracting,
and validating evidence relating to UK higher education frameworks, levels,
awards, regulatory oversight, qualification comparability, and transnational
education (TNE) delivery with Sri Lankan partners.

\subsection{Source Categories}

Two categories of sources were defined:

\begin{itemize}
    \item \textbf{Fixed UK sources:} Office for Students (OfS); Quality Assurance 
    Agency for Higher Education (QAA); Department for Education (DfE); UK~ENIC 
    (ECCTIS); Universities~UK; SCQF Partnership; Scottish Funding Council (SFC); 
    Welsh Government / CTER; Department for the Economy Northern Ireland (DfE~NI); 
    Ofqual.

    \item \textbf{Context sources (Sri Lanka):} Sri Lanka Qualifications Framework 
    (SLQF) / University Grants Commission (UGC); British Council Sri~Lanka.
\end{itemize}

\noindent\textit{Note.} Potential Sri Lankan sources identified during protocol development 
(MOHE, TVEC) yielded no eligible documents under the inclusion criteria and were 
therefore not included in the final evidence set.



Only official, authoritative, and versioned documents from these providers were
eligible for inclusion.

\subsection{Time Window}

The review prioritised current versions in force as of~2025. Earlier versions
were consulted for lineage or regulatory transition where required, generally
back to approximately~2008.

\subsection{Search Strategy}

Eight predefined query groups (G1--G8) were executed verbatim on Google and on
provider-domain search bars. Searches combined:

\begin{itemize}
    \item site-restricted queries;
    \item award/credit/level terminology;
    \item regulatory terms (conditions, quality, validation, franchise, TNE);
    \item recognition/comparability terminology (ENIC, GCE~A/L, Sri~Lanka);
    \item exit-award and intermediate-award terms;
    \item partnership/oversight terms.
\end{itemize}

All queries were executed using identical syntax for reproducibility.

\subsection{Inclusion Criteria}

Items were included where they met one or more of the following:

\begin{itemize}
    \item official frameworks, award descriptors, qualification statements;
    \item regulatory frameworks, conditions, sector-recognised principles;
    \item official guidance for validation, franchise, or TNE;
    \item university or awarding-body academic regulations (exit awards);
    \item authoritative sector reports or peer-reviewed analyses.
\end{itemize}

\subsection{Exclusion Criteria}

Items were excluded where they met one or more of the following:

\begin{itemize}
    \item opinion/marketing/blog/agent materials;
    \item non-UK or non-HE materials;
    \item superseded versions without lineage relevance;
    \item general information without regulatory or standards substance;
    \item documents without year, authority, or identifiable provenance.
\end{itemize}

\subsection{DocID Scheme}

Each screened item was assigned a DocID of the form 
\texttt{SRC-YY-\#\#\#} (where SRC is replaced by the source-provider prefix, 
e.g., QAA, OFS, SLQF, BCSL). This identifier linked the screening log, evidence 
table, PDF archive, and reference list.


\subsection{Screening Process}

Screening was logged in \texttt{screening\_log\_\_v0.1\_freeze\_2025-11-24.csv} with fields:

\begin{itemize}
    \item DocID;
    \item Source;
    \item URL and capture details;
    \item Title;
    \item Publisher;
    \item Year;
    \item Country;
    \item DocType;
    \item Status (Include / Exclude / Borderline);
    \item ReasonCode (I1--I4, E1--E4, B1);
    \item Notes;
\end{itemize}

\raggedright
PDFs were saved under a consistent directory scheme. For example: \texttt{evidence/QAA/QAA-24-001\_\_QAA\_\_2024.pdf}. These files were logged in \texttt{evidence\_manifest\_\_v0.1\_freeze\_2025-11-24.csv}.


\subsection{Second-Reader Sampling}

A second reader independently reviewed:

\begin{itemize}
    \item approximately 15\% of included items,
    \item all Borderline items,
    \item all Excluded items.
\end{itemize}

Disagreements were resolved by discussion and summarised in
\texttt{agreement\_summary\_\_v0.1\_2025-11-24.csv}.

\subsection{Data Extraction}

Extraction followed the schema in \texttt{evidence\_table\_\_v0.1\_freeze\_2025-11-24.csv}, including:

\begin{itemize}
    \item DocID, Title, Year, Jurisdiction, Document Type;
    \item Primary Focus;
    \item Key Provisions / Level Descriptors;
    \item Levels and Credits;
    \item Exit Awards;
    \item TNE or Partnership Rules;
    \item Recognition / Entry (ENIC);
    \item Citations and Use in Review.
\end{itemize}

\subsection{Risk and Limitations}

Risks included:

\begin{itemize}
    \item regulatory updates during the review period;
    \item site restructures affecting URL access;
    \item unpaginated PDFs;
    \item sector reports that are descriptive rather than prescriptive.
\end{itemize}

\subsection{Deviations}

No deviations occurred during the review.

\subsection{ReasonCode Dictionary}

The following ReasonCodes were used during screening and second-reader
verification to ensure consistency, reproducibility, and auditability across all
decisions. Codes are grouped into Inclusion (I1--I4), Exclusion (E1--E4), and
Borderline (B1).

\begin{table}[h!]
\centering
\begin{tabular}{p{2cm} p{11cm}}
\toprule
\textbf{Code} & \textbf{Definition} \\
\midrule
I1 & Meets inclusion criteria: official framework, qualification descriptor, or award statement. \\
I2 & Meets inclusion criteria: regulatory framework, condition, or sector-recognised principle. \\
I3 & Meets inclusion criteria: official validation, franchise, or TNE guidance. \\
I4 & Meets inclusion criteria: authoritative sector report or peer-reviewed analysis. \\
\midrule
E1 & Excluded: opinion, marketing, agent, promotional, or blog content. \\
E2 & Excluded: non-UK or not higher-education-level material. \\
E3 & Excluded: superseded or outdated version without lineage relevance. \\
E4 & Excluded: insufficient authority, missing year, unverifiable provenance, or unclear source. \\
\midrule
B1 & Borderline relevance: requires additional judgement, context interpretation, or second-reader review. \\
\bottomrule
\end{tabular}
\end{table}
